\section{Related work}
Kohavi and John evaluated feature subsets by training the intended predictor as part of a wrapper~\cite{kohavi1997wrappers}. Recursive feature elimination repeatedly fits the estimator as variables are removed~\cite{guyon2002rfe}. Both approaches can be costly when many candidate subsets are examined. A masked prediction is cheaper once the original model exists. Carter et al.\ used backward masking to find inputs sufficient for a single prediction~\cite{carter2019sufficient}. Here deletion applies to whole columns and the quantity of interest concerns a new fit. Covert et al.\ distinguished several removal operations and their estimands~\cite{covert2021removing}. Hooker et al.\ evaluated rankings by removing features and retraining~\cite{hooker2019roar}. Our proposal-level comparison follows that distinction by holding the proposed set fixed while changing the evaluator.

MinShap takes the smallest conditional feature importance across conditioning sets in a characterization of redundancy~\cite{zheng2026minshap}. ShapG samples coalitions informed by a graph of feature dependencies~\cite{zhao2025shapg}. Their unit of evaluation is a feature's attributed contribution. In our deletion test the worst context can veto a block, and the decision depends on the fitted predictor, metric and sampled rows. The block search shares subdivision and a single-deletion stopping rule with delta debugging~\cite{zeller2002ddmin}. We use these existing ideas to study disagreement between decisions under masking and decisions under refitting.

Benchmark comparisons depend on the evaluation criterion and the learner applied after feature selection~\cite{bommert2020filters}. Selection on noisy validation scores also introduces an opportunity to overfit a model choice~\cite{cawley2010selection}. We therefore preserve outer test rows during each selector fit and identify analyses that revisit the same task scores after initial evaluation.
